\documentclass[10pt,a4paper]{amsart}
\usepackage[margin=19mm]{geometry}
\usepackage{amsmath,amssymb,mathtools}
\usepackage[scr=rsfs,cal=euler]{mathalfa}
\usepackage{xfrac}
\usepackage[unicode,hidelinks,bookmarks=false]{hyperref}
\newcommand{\ie}{i.e.\ }
\newcommand{\cf}{cf.\ }
\newcommand{\resp}{respectively\ }
\newcommand{\Subsection}{\subsection}
\providecommand{\nobreakdash}{\nobreak\hbox{-}}
\setlength{\emergencystretch}{2em}
\multlinegap0pt
\usepackage{amssymb}
\usepackage{dsfont}
\usepackage[shortlabels]{enumitem}
\usepackage{tikz}
\newtheorem{theorem}{Theorem}
\def\BK{{\bf K}}
\def\bLa{{\bf \Lambda}}
\def\bS5{{\bf S5}}
\def\mo{\vDash}
\def\rtt{\rightthreetimes}
\def\sbq{\subseteq}
\def\sqr{\square}
\def\ti{\times}
\title{Some prospects for semiproducts and products of modal logics}
\author{Valentin Shehtman, Dmitry Shkatov}
\date{Source: arXiv:2607.17928v1 (CC BY 4.0)}
\begin{document}
\maketitle

\noindent\emph{Source and licence.} Some prospects for semiproducts and products of modal logics. Version arXiv:2607.17928v1 (CC BY 4.0), distributed by arXiv under the Creative Commons Attribution 4.0 International licence, http://creativecommons.org/licenses/by/4.0/, as recorded in the arXiv licence metadata for this exact version and shown on the abstract page https://arxiv.org/abs/2607.17928v1. Full licence notice: \texttt{licenses/arxiv-2607.17928-CC-BY-4.0.txt}.\par\smallskip

\noindent\textit{Editorial scope.} Counted unit: the article's theorem Tneg (source0.tex lines 919--923). The article's complete original proof of it is delivered with it (source0.tex lines 925--935, one \texttt{proof} environment of 11 lines). No mathematics is written, repaired or re-derived: the statement and the proof are the article's own lines, in the article's own notation, and no retained line is substituted or re-flowed.  No further source-local context is needed: every cross-reference of every retained line resolves inside the two delivered spans.  Nothing was omitted from any retained span, so the package carries no omission record.  No retained line contains a citation, so the wrapper carries no bibliography and no bibliography entry.  No retained line contains a figure environment, an image-inclusion command or drawing code, so no image is delivered and the package declares no asset dependency.  Editorial formatting only, in addition: an AMS \texttt{amsart} preamble that restates the article's own declarations and macros used by the retained lines, namely the environments \texttt{theorem} on separate counters, together with the standard packages those lines need.  The retained environments print in this document as Theorem 1 in order of appearance, whereas the article prints the same items with the numbering of its own full document.  The complete native source of the article as distributed in the arXiv e-print for this exact version is bundled unchanged as \texttt{sources/205612/source0.tex}.
\begin{theorem}\label{Tneg}
  If $\BK+Alt_n\sbq\bLa\sbq\BK+Alt_n+\sqr^m\bot$ for $n\geq 3$,
  $m\geq 2$, then
  $\bLa$ and $\bS5$ are neither product- nor semiproduct-matching.
\end{theorem}

\begin{proof} (Sketch.) Take the product $F_1\ti F_2$, where $F_1$ is
  the irreflexive tree with the root $0$ and the leaves $1,\ldots,n$
  and $F_2$ is the two-element cluster $\{1,2\}$; replace $R_v$ by the
  least equivalence relation $S_2$ such $(x,y)S_2(x',y')$ for $x=x'=0$
  or $x=x'>3$,
  $(1,1)S_2(2,2), ~(1,1)S_2(3,2),~(1,2)S_2(2,1), ~(1,2)S_2(3,1)$. The
  resulting frame $G_n$ is not a p-morphic image of a semiproduct of a
  $(\BK+Alt_n)$-frame and a cluster while
  $G_n\mo [\BK+Alt_n+\sqr^2\bot,\bS5]$. Therefore its Fine-Jankov
  formula belongs to $(\bLa\rtt \bS5)-[\bLa,\bS5]$.
\end{proof}

\end{document}
